%% Hand-authored circuitikz figure -- NOT generated by maestro2tex.
%% Topology transcribed device-by-device from the netlist of
%%   .simulation/castalia/tb_anatop_pixel_dacR2R12/maestro/results/maestro/
%%   Interactive.34/1/DacNonlinTB/netlist/input.scs
%% (subckt anatop_pixel_dacR2R12, 76 unit resistors + 12 AND2X8MA10TH drivers).
%% Device dimensions are proprietary: only unit counts and ratios appear here.
\begin{figure}[htbp]
  \centering
  \providecommand{\dacbv}[1]{\ensuremath{\langle #1\rangle}}
  \resizebox{\linewidth}{!}{%
\ctikzset{bipoles/length=0.9cm, resistors/scale=0.85, logic ports=american, logic ports/scale=1.05}
\begin{circuitikz}[
    every node/.append style={font=\footnotesize},
    nlab/.style={font=\scriptsize, inner sep=1.5pt},
    gatelab/.style={font=\scriptsize, inner sep=1pt},
    note/.style={font=\scriptsize, align=left, inner sep=2pt},
    railnode/.style={circle, fill=black, inner sep=1.05pt},
  ]

% ---- geometry -------------------------------------------------------
\def\xT{0}    \def\xa{2.6}  \def\xb{5.2}
\def\xgl{7.0} \def\xgr{8.6} \def\xj{10.4} \def\xk{13.0} \def\xout{14.6}
\def\yLad{5.0} \def\yGo{2.9} \def\yEn{0.55} \def\yA{-0.35}

% ===== ladder rail ===================================================
\draw (\xT,\yLad) -- (\xa,\yLad);
\draw (\xa,\yLad) to[R, l=$R$] (\xb,\yLad);
\draw (\xb,\yLad) to[R, l=$R$] (\xgl,\yLad);
\draw[dashed] (\xgl,\yLad) -- (\xgr,\yLad);
\draw (\xgr,\yLad) to[R, l=$R$] (\xj,\yLad);
\draw (\xj,\yLad) to[R, l=$R$] (\xk,\yLad);
\draw (\xk,\yLad) -- (\xout,\yLad);
\node[railnode] at (\xout,\yLad){};
\node[nlab, right=2pt] at (\xout,\yLad){\texttt{Out}};

% ---- tap nodes ------------------------------------------------------
\foreach \x in {\xa,\xb,\xj,\xk} { \node[railnode] at (\x,\yLad){}; }
\node[nlab, above=3pt] at (\xa,\yLad){\texttt{n}\dacbv{1}};
\node[nlab, above=3pt] at (\xb,\yLad){\texttt{n}\dacbv{2}};
\node[nlab, above=3pt] at (\xj,\yLad){\texttt{n}\dacbv{11}};

% ---- elision label --------------------------------------------------
\node[nlab, align=center, above=17pt] at ({0.5*(\xgl+\xgr)},\yLad)
  {bits 2--9:\\8 identical\\$R$ and $2R$ sections};
\node[font=\large] at ({0.5*(\xgl+\xgr)},\yGo+0.9){$\cdots$};

% ===== terminating 2R leg ============================================
\draw (\xT,\yLad) -- (\xT,4.5);
\draw (\xT,4.5) to[R, l_=$R$] (\xT,3.1);
\node[railnode] at (\xT,3.1){};
\node[nlab, right=3pt] at (\xT,3.1){\texttt{n}\dacbv{0}};
\draw (\xT,3.1) to[R, l_=$R$] (\xT,1.7);
\draw (\xT,1.7) -- (\xT,1.15) node[ground, scale=0.8]{};
\node[railnode] at (\xT,\yLad){};
\node[note, anchor=north east] at (\xT-0.55,4.3)
  {terminating $2R$:\\rung\,0 $+$ \texttt{RTERM1,2}\\(no bit taps here)};

% ===== En rail (drawn first; A leads cut a gap through it) ===========
\draw[thick] (\xa-1.5,\yEn) -- (\xgl,\yEn);
\draw[thick, dashed] (\xgl,\yEn) -- (\xgr,\yEn);
\draw[thick] (\xgr,\yEn) -- (\xk+0.8,\yEn);
\node[nlab, left=2pt] at (\xa-1.5,\yEn){\texttt{En}};

% ===== bit slices ====================================================
% #1 x, #2 bit index
\newcommand{\bitslice}[2]{%
  \draw (#1,\yLad) -- (#1,4.5);
  \draw (#1,4.5) to[R, l_=$2R$] (#1,\yGo);
  \node[and port, rotate=90] (G#2) at (#1,\yGo-1.05) {};
  \draw (G#2.out) -- (#1,\yGo);
  \node[railnode] at (#1,\yGo){};
  \node[nlab, right=2pt] at (#1,\yGo){\texttt{f}\dacbv{#2}};
  \draw (G#2.in 1) -- (G#2.in 1 |- 0,\yEn);
  \node[railnode] at (G#2.in 1 |- 0,\yEn){};
  \draw[white, line width=2.6pt] (G#2.in 2) -- (G#2.in 2 |- 0,\yA);
  \draw (G#2.in 2) -- (G#2.in 2 |- 0,\yA);
  \node[nlab, below=1pt] at (G#2.in 2 |- 0,\yA){\texttt{A}\dacbv{#2}};
}
\bitslice{\xa}{0}
\bitslice{\xb}{1}
\bitslice{\xj}{10}
\bitslice{\xk}{11}

% ===== annotations ===================================================
\node[note, anchor=north west, draw, rounded corners, inner sep=4pt] at (\xT-1.9,\yLad+2.6)
  {12 bits, one $2R$ branch each.\\
   Every resistor is one \emph{unit} device:\\
   $R=2$ units (\texttt{RS1}$+$\texttt{RS2}),\\
   $2R=4$ units (\texttt{RP1}--\texttt{RP4})};
\node[note, anchor=west] at (\xk+1.1,\yGo-1.0)
  {\texttt{AND2X8MA10TH}\\bit drivers, supplied\\from \texttt{inh\_vdd}};

\end{circuitikz}
  }
  \caption[{The 12-bit voltage-mode R-2R ladder.}]{The 12-bit voltage-mode R--2R ladder of \texttt{anatop\_pixel\_dacR2R12}, transcribed from the \texttt{DacNonlinTB} netlist. Bit $k$ drives tap \texttt{n}$\langle k{+}1\rangle$ through its own $2R$ branch, so the twelve taps are \texttt{n}$\langle 1\rangle$--\texttt{n}$\langle 11\rangle$ plus \texttt{Out}, and the MSB branch lands directly on \texttt{Out} at the top of the ladder; \texttt{n}$\langle 0\rangle$ carries no bit, its rung and \texttt{RTERM1,2} together forming the terminating $2R$. Bits 2--9 are elided: eight further sections identical to those drawn sit between \texttt{n}$\langle 2\rangle$ and \texttt{n}$\langle 11\rangle$. Every element is the same unit resistor, so the network itself is exact and gives $\mathrm{Out}=(\mathrm{code}/4096)\,V_{\mathrm{ref}}$; the reference is the \emph{output high level} of the \texttt{AND2X8MA10TH} drivers, which sit in series with each $2R$ branch and bound the ladder in place of the rails --- measured zero scale is \SI{90}{\micro\volt} above $V_{\mathrm{SS}}$ and full scale \SI{103}{\micro\volt} below the ideal $(4095/4096)\,V_{\mathrm{ref}}$ on a \SI{2.5}{\volt} rail. Device dimensions are proprietary and are not stated.}
  \label{fig:dacr2r12-ladder}
\end{figure}
